\documentclass{article}
\usepackage{epcc}
\usepackage{hyperref}
\usepackage{listings}


\title{Fractals, Task Farms and Load Imbalance: The Outcomes}
\date{}


\begin{document}

\makeEPCCtitle

\vspace*{-1cm}


This is a summary of what should be learned in the Fractals exercise. At the bottom is a set of results which can help with some of the questions posed in the exercise.

\section{What We Should Have Accomplished}

In the Fractals exercise, the following should have been attempted:
\begin{itemize}
\item Accessing a remote HPC resource.
\item Obtained and compiled source code.
\item Assessed the performance of a task farm across configurations.
\end{itemize}

\section{What We Should Take Away}

After attempting the exercise, it is important to reinforce a few key points. 
\textbf{Note:} if anything is unclear or missing, ask the
demonstrators who are there to help.

\subsection{Key Points}


\textbf{Task Farm}, also known as the master/worker pattern, allows a
master process to distribute work to a set of workers processes. This
is suitable for calculations which can be split up into smaller chunks
with little or no inter-worker communications. The pattern can be used
for other types of tasks but it complicates the situation and other
patterns may be more suitable for implementing. The master process is
responsible for creating, distributing and gather the individual
jobs. This will often include the collation of the results into an
interpretable whole.


\textbf{Tasks} are the units of work in a task farm. These can vary in
size and do not have to be of consistent execution time. If execution
times are known it can help with load balancing. For a task farm, the
tasks should as far as possible be discrete units of work from the
overall calculation.


\textbf{Queues} are what the pool of tasks a master generates can be
thought of. Each worker will be given a new task when they finish
their current task until the pool is empty. If the scheme is to give
out tasks as and when requested, then the pool will act like a
queue. If there is some filtering of tasks, like known long and short
tasks, they may be separated into multiple queues. Some systems will
generate the tasks on-the-fly while others will prepare the pool of
tasks before starting to distribute tasks and others will take a
strategy between these two.


\textbf{Load Balancing} is how a system determines how work or tasks
are distributed across processing elements to ensure good
performance. Successful load balancing will avoid overloading a single
element, maximising the throughput of the system and make best use of
resources available. Poor load balancing will lead to elements of the
system being under-utilised and reducing the performance as it gets
dominated by over-subscribed resources.


The distribution of tasks is generally done in either a pre-planned or
responsive mode. In responsive mode, the master will give out tasks as
requested by workers. In pre-planned mode, where task lengths are
known or estimated, the master will map out a set of tasks for each
worker to complete. This allows a clear picture of the overall
workload to be developed: most planned systems  allocate larger
tasks first.

\subsection{Other}

\textbf{Cost} is increasingly becoming a big factor in how resources
are requested and used. Hopefully this exercise will have highlighted
that often just throwing more power and resources at something
eventually has diminishing returns in performance. This means the
monetary cost can be a lot higher for only minor returns over a more
efficient and less expensive configuration at run time. Is it
necessary to run on 4000 cores, or can the same problem be run on 2000
cores?


\textbf{Generic patterns and techniques} are common to software
development and the task farm is a common pattern to recognise. This
is commonly use in parallel and high performance computing and while
the exercise focus on a fractal problem the underlying framework could
be adapted to any problem with similar properties.

\section{Sample Results}


\subsection{Cores, Time, Worker Workloads}
Example results for the default image size (768 $\times$ 768 pixels), fixed number of iterations (5000), 64 tasks of the size of 96 x 96 pixels and varying number of workers :

\begin{table}[h]
\begin{center}
\begin{tabular}{ | c | c | c | c | c | c | }
\hline
	Workers & Time (s) & Avg. Workload & Max Workload & Min Workload & Load Imbalance Factor \\ \hline
	4 & 1.52 &124505763 & 126124383 & 122696852 & 1.01 \\ \hline
	8 & 0.96 &  62252881 & 77744803 & 51117022  & 1.25 \\ \hline
	16 & 0.59 & 31126440 & 46737752 & 10968369  & 1.50 \\ \hline
	32 & 0.57 & 15563220  &  46114456  & 67614 & 2.96 \\ \hline
	64 & 0.57 & 7781610  & 46089216&  9246 &5.92 \\ \hline	\end{tabular}
	 \caption{Example Run-Times for different number of workers and their Avg/Min/Max Workloads.}
	 \end{center}
	 \end{table}


\subsection{Time, Task Size, Load Imbalance  }
Example results for the default image size (768 $\times$ 768 pixels), fixed number of iterations (5000), fixed number of workers (16) and varying number of tasks :

\begin{table}[h]
\begin{center}
\begin{tabular}{ | l | l | l | l | l | }
\hline
	Number of Tasks (Task Size)  & Time (s) & Load Imbalance Factor \  \\ \hline
        16 (192 $\times $ 192 ) & 1.93 & 5.034  \\ \hline
	64 (96 $\times$ 96 ) & 0.59 & 1.501 \\ \hline
	256 (48 $\times$ 48) & 0.43 & 1.108  \\ \hline
	4096 (12 $\times$ 12) & 0.4 & 1.017 \\ \hline
	36864 (4 $\times$ 4) & 0.4 & 1.003  \\ \hline
	147456 (2 $\times$ 2) & 0.47 & 1.017 \\ \hline
	589824 (1 $\times$ 1) & 0.80 & 1.006 \\ \hline
        \end{tabular}
        \caption{Example execution Times for 16 workers and varying number of Tasks.}
	 \end{center}
	 \end{table}
% \begin{table}[h]
% \begin{center}
% \begin{tabular}{ | l | l | l | l | l | }
% \hline
% 	Calculation Time  &  (seconds) & \  & \  & \  \\ \hline
%          Cores / Image Size & 1000 & 2000 & 4000 & 8000 \\ \hline
% 	4 & 2.19 & 8.75 & 35.04 & 140.05\\ \hline
% 	8 & 1.00 & 3.98 & 15.92 & 63.66 \\ \hline
% 	16 & 0.49 & 1.91 & 7.50 & 29.97 \\ \hline
% 	32 & 0.25 & 0.97 & 3.67 & 14.51 \\ \hline
% \end{tabular}
%  \caption{Example Calculation Times for Core and Image Sizes}
%  \end{center}
%  \end{table}

\end{document}

